% !TeX program = xelatex
% Faithful transcription of eighteen handwritten source pages, in image order.
% Medical statements and English spellings are preserved as written.
\ifdefined\settextfont
  \def\GIAStandalone{0}
\else
  \def\GIAStandalone{1}
  \documentclass[11pt,a4paper]{article}
  \usepackage[margin=1.8cm]{geometry}
  \usepackage{amsmath,amssymb,enumitem}
  \usepackage{xepersian}
  \IfFontExistsTF{Amiri}{\settextfont{Amiri}}{\settextfont{FreeSerif}}
  \setlatintextfont{FreeSerif}
  \setlength{\parindent}{0pt}
\fi
\providecommand{\en}[1]{\lr{#1}}
\newcommand{\GIAarrow}{\en{$\longleftarrow$}}
\newcommand{\GIAheader}{{\footnotesize\en{Note book\qquad Subject: 18\qquad Veor: \_\_\_\quad Month: \_\_\_\quad Date: \_\_\_}\par}}
\newcommand{\GIAfooter}[1]{\par\medskip{\footnotesize\en{BAHAR DANESH ASL}\hfill #1\par}}
\ifnum\GIAStandalone=1 \begin{document}\fi

% pic01.jpg = 20261019_082636808.jpg
\GIAheader
\section*{\en{Gastrointestinal, Hepatobiliary, urinary tracked abnormalities}}
\begin{enumerate}[label=\arabic*.,leftmargin=*,itemsep=.5em]
\item اغلب ازوفاگوگرام‌ها به‌صورت بای‌فازیک انجام می‌شوند.
\begin{latin}Single-contrast + Double-contrast\end{latin}
\item \textbf{تصویربرداری انتخابی برای تعیین آسپیراسیون جسم خارجی}
\GIAarrow\ \textbf{\en{Video esophagography}}
\begin{itemize}[nosep]\item برای بررسی مکانیسم بلع استفاده می‌شود.\end{itemize}
\item فلوروسکوپی مری \GIAarrow\ اختلالات حرکتی مری را نشان می‌دهد.
\item \textbf{باریوم ازوفاگوگرام تصویربرداری اولیه به \GIAarrow\ دیسفاژی}
\item \en{Appearance of esophageal carcinoma}
\begin{latin}\begin{itemize}[nosep]
\item Annular constricting
\item Tumor encircles the normal lumen and almost completely obstructs.
\item Polypoid mass
\item Irregular and rigid
\end{itemize}\end{latin}
\end{enumerate}
\subsection*{هیاتال هرنیا و \en{GERD}}
\begin{enumerate}[label=\arabic*.,leftmargin=*]
\item انواع هیاتال هرنیا
\begin{itemize}[nosep]
\item \en{Sliding}: \en{GEJ} بالاتر از دیافراگم است.
\item پاراازوفاژیال: \en{GEJ} پایین‌تر از دیافراگم است و بخشی از معده بالاتر.
\end{itemize}
\end{enumerate}
\GIAfooter{۱}

% pic02.jpg = 20261019_082641908.jpg
\clearpage\GIAheader
\begin{enumerate}[label=\arabic*.,leftmargin=*,itemsep=.5em]\setcounter{enumi}{1}
\item به‌طور کلی شیوع هیاتال هرنیا با افزایش سن افزایش می‌یابد.
\item اغلب هیاتال هرنیاها بدون علامت هستند.
\item \en{GERD} حتی به افراد بدون هیاتال هرنیا هم دیده می‌شود.
\begin{itemize}[nosep]\item معمولاً به علت اختلال عملکرد اسفنکتر تحتانی مری\end{itemize}
\item یافته‌های تصویربرداری هیاتال هرنیا
\begin{latin}\begin{itemize}[itemsep=.4em]
\item A bulbous area of the distal esophagus containing oral contrast at the level of the diaphragm, with failure of esophagus to narrow as it passes through the esophageal hiatus.
\item Extension of multiple gastric folds above the diaphragm
\item Visualization of a thin, circumferential filling defect in the distal esophagus $\longrightarrow$ \textbf{Schatzki Ring}
\end{itemize}\end{latin}
\item در فلوروسکوپی دیدن ریفلاکس ضرورتاً به معنی وجود علائم \en{GERD} نیست و ندیدن آن هم ردکنندهٔ \en{GERD} نخواهد بود.
\end{enumerate}
\subsection*{معده و دئودنوم}
\begin{enumerate}[label=\arabic*.,leftmargin=*]
\item امروزه برای بررسی معده
\begin{itemize}[nosep]
\item لومن: اندوسکوپی فوقانی
\item دیواره و ساختمان: \textbf{\en{CT} با کنتراست}
\end{itemize}
\end{enumerate}
\GIAfooter{۲}

% pic03.jpg = 20261019_082650112.jpg
\clearpage\GIAheader
\begin{enumerate}[label=\arabic*.,leftmargin=*,itemsep=.6em]\setcounter{enumi}{1}
\item زخم معده
\begin{itemize}[itemsep=.3em]
\item هلیکوباکتر عامل \en{$\frac{3}{4}$} موارد زخم معده است.
\item \en{NSAID} عامل سایر موارد.
\item شایع‌ترین محل‌های زخم \GIAarrow\ لسرکورویچر و پره‌پیلوریک و بادی و آنتر معده
\item ۹۵\% زخم‌های معده خوش‌خیم هستند.
\end{itemize}
\item کارسینوم معده
\begin{itemize}[itemsep=.3em]
\item اغلب آن‌ها در \en{$\frac{1}{3}$} تحتانی معده و در لسرکورویچر رخ می‌دهند.
\item تصویربرداری \GIAarrow\ \en{CT} شکم؛ \en{UGI} دابل کنتراست
\item از \en{CT} برای استیجینگ هم استفاده می‌شود.
\item انواع فرم کارسینوم معده: \en{polypoid}؛ \en{ulcerative}؛ \en{Infiltrating (linitis plastica)}
\item توده‌هایی که شبیه کارسینوم هستند \GIAarrow\ لایومیوم؛ لنفوم
\end{itemize}
\item زخم دئودنوم
\begin{itemize}[itemsep=.3em]
\item ۲--۳ برابر شایع‌تر از زخم معده هستند.
\item تقریباً همهٔ آن‌ها در بلب دئودنوم هستند.
\item اکثریت در دیوارهٔ خلفی
\item ۹۵\% ناشی از هلیکوباکتر
\end{itemize}
\end{enumerate}
\GIAfooter{۳}

% pic04.jpg = 20261019_082655585.jpg
\clearpage\GIAheader
\begin{itemize}[itemsep=.4em]
\item \en{UGI} دابل کنتراست حساسیت ۹۰\% در تشخیص زخم دئودنوم
\item \textbf{بهترین ابزار تشخیص عوارض زخم دئودنوم \GIAarrow\ \en{CT}}
\begin{latin}Obstruction\\Perforation\\Penetration\\Hemorrhage\end{latin}
\end{itemize}
\subsection*{روده بزرگ و کوچک}
\begin{enumerate}[label=\arabic*.,leftmargin=*,itemsep=.5em]
\item مستقل از تکنیک تصویربرداری برای بررسی روده استفاده می‌کنیم، لومن روده باید \en{distention} و \en{opacification} داشته باشد تا بررسی دقیق باشد.
\begin{itemize}[nosep]
\item برای انجام \en{CT}، کنتراست خوراکی به دوزهای منقسم داده می‌شود.
\item بجز در \en{CT} برای تروما و سنگ و عروق
\end{itemize}
\item بیماری کرون
\begin{itemize}[itemsep=.3em]
\item یک بیماری مزمن، عودکننده، گرانولوماتوز و التهابی است.
\item به‌طور شایع ایلئوم و کولون راست را درگیر می‌کند.
\item \en{Skip area} دیده می‌شود.
\item مستعد تشکیل فیستول است.
\item تصویربرداری \GIAarrow\ \textbf{سی‌تی شکم و لگن}؛ \en{Barium small bowel follow-through}
\end{itemize}
\end{enumerate}\GIAfooter{۴}

% pic05.jpg = 20261019_082702650.jpg
\clearpage\GIAheader
\begin{enumerate}[label=\arabic*.,leftmargin=*,itemsep=.5em]\setcounter{enumi}{2}
\item یافته‌های تصویربرداری در کرون
\begin{latin}\begin{itemize}[itemsep=.5em]
\item Narrowing, Irregularity and ulceration of the terminal ileum with proximal small bowel dilation.
\item Separation of the loops of bowel due to fatty infiltration of the mesentery (Proud loop).
\item \textbf{String sign}, narrowing of the terminal ileum into a near slit-like opening by spasm and fibrosis.
\item \textbf{Fistulae, specially between the ileum and colon}
\end{itemize}\end{latin}
\item دایورتیکولوز
\begin{itemize}[itemsep=.3em]
\item دایورتیکول‌ها حاوی هوا هستند.
\item شیوع آن‌ها با افزایش سن بالا می‌رود.
\item معمولاً متعدد هستند.
\item در ۹۰\% موارد بی‌علامت هستند.
\item دایورتیکولوز شایع‌ترین علت مسیو \en{lower GI bleeding} هستند.
\item زمانی که خونریزی رخ دهد دایورتیکول‌های سمت راست بیشتر از سمت چپ خونریزی می‌کنند.
\item بیشتر در کولون سیگموئید دیده می‌شوند.
\item تصویربرداری \GIAarrow\ باریوم انما؛ \en{CT}
\end{itemize}
\end{enumerate}\GIAfooter{۵}

% pic06.jpg = 20261019_082708752.jpg
\clearpage\GIAheader
\begin{enumerate}[label=\arabic*.,leftmargin=*,itemsep=.5em]\setcounter{enumi}{4}
\item دایورتیکولیت
\begin{itemize}[itemsep=.3em]
\item شایع‌ترین علت آن انسداد یا تحریک مخاطی دایورتیکول است.
\item تصویربرداری انتخابی \GIAarrow\ \en{CT}
\item یافته‌های تصویربرداری دایورتیکولیت در \en{CT}
\item دیدن دایورتیکول، علاوه موارد زیر:
\end{itemize}
\begin{latin}\begin{itemize}[nosep]
\item \textbf{Thickening of the adjacent colonic wall ($>4$ mm)}
\item Pericolonic inflammation
\item Abscess formation
\item Perforation of the colon
\end{itemize}\end{latin}
\item پولیپ‌های کولون
\begin{itemize}[itemsep=.3em]
\item ریسک بدخیمی با سایز پولیپ افزایش می‌یابد.
\item شایع‌ترین پولیپ‌های کولون \GIAarrow\ پولیپ هایپرپلاستیک که ریسک بدخیمی ندارند.
\item پولیپ‌های آدنوماتوز بالای \en{1.5 cm} حدود ۱۰\% ریسک بدخیمی دارند.
\item انواع پولیپ \GIAarrow\ \en{Sessile}؛ پدانکوله
\end{itemize}
\item یافته‌های تصویربرداری کارسینوم کولون
\begin{latin}\begin{itemize}[nosep]
\item Persistent, large polypoid filling defect.
\item Annular constriction of the lumen: \textbf{Apple-core lesion}
\end{itemize}\end{latin}
\end{enumerate}\GIAfooter{۶}

% pic07.jpg = 20261019_082715731.jpg
\clearpage\GIAheader
\begin{latin}\begin{itemize}[nosep]
\item Frank perforation or microperforation
\item Large bowel obstruction
\item Metastasis $\longrightarrow$ liver and lung
\end{itemize}\end{latin}
\begin{enumerate}[label=\arabic*.,leftmargin=*,itemsep=.5em]\setcounter{enumi}{7}
\item کولیت
\begin{itemize}[nosep]
\item علل: عفونی؛ اولسراتیو و گرانولوماتوز؛ ایسکمیک؛ ناشی از پرتودرمانی؛ مرتبط با آنتی‌بیوتیک
\item یافته‌های تصویربرداری کولیت در \en{CT}
\end{itemize}
\begin{latin}\begin{itemize}[nosep]
\item Segmental thickening of the bowel wall
\item Irregular narrowing of the bowel lumen due to edema (Thumb-printing)
\item Infiltration of the surrounding fat
\item Dual pattern (Accordion sign)
\end{itemize}\end{latin}
\item آپاندیسیت
\begin{itemize}[itemsep=.3em]
\item \textbf{\en{CT} و \en{US} تصویربرداری‌های تشخیصی انتخابی هستند.}
\item از \en{MRI} هم می‌توان استفاده نمود.
\item آپاندیکولیت: کلسیفیه شدن آپاندیس در ۱۰\% افراد
\item در \en{CT} خیلی خوب دیده می‌شود.
\item درد شکم + آپاندیکولیت \GIAarrow\ در ۹۰\% موارد آپاندیسیت است.
\item ریسک بالاتر پرفوریشن
\end{itemize}
\end{enumerate}\GIAfooter{۷}

% pic08.jpg = 20261019_082721274.jpg
\clearpage\GIAheader
\begin{enumerate}[label=\arabic*.,leftmargin=*]\setcounter{enumi}{9}
\item یافته‌های تصویربرداری آپاندیسیت حاد در \en{CT}
\begin{latin}\begin{itemize}[nosep]
\item Dilated appendix ($>6$ mm) that does not fill with oral contrast.
\item Periappendiceal inflammation
\item Increased contrast enhancement of the wall of the appendix due to inflammation
\item Perforation (in up to 30\% of cases)
\end{itemize}
\textbf{* The presence of a large amount of free intraperitoneal air should point to another diagnosis.}
\end{latin}
\end{enumerate}
\subsection*{خونریزی گوارشی تحتانی \en{(LGIB)}}
\begin{enumerate}[label=\arabic*.,leftmargin=*,itemsep=.5em]
\item تصویربرداری‌هایی که در \en{LGIB} استفاده می‌شوند:
\begin{itemize}[nosep]
\item کولونوسکوپی
\item کاتتر آنژیوگرافی
\item \en{CT} آنژیوگرافی \en{(CTA)}
\item \en{RNBS}
\end{itemize}
\item اگر همودینامیک بیمار پایدار باشد انتخاب ما کولونوسکوپی است.
\item اگر بیمار ناپایدار باشد انتخاب ما کاتتر آنژیوگرافی می‌باشد.
\item در \en{CTA} باید حداقل \en{0.5 ml/min} خونریزی وجود داشته باشد تا نشان بدهد.
\end{enumerate}\GIAfooter{۸}

% pic09.jpg = 20261019_082727801.jpg
\clearpage\GIAheader
\subsection*{پانکراس}
\begin{enumerate}[label=\arabic*.,leftmargin=*,itemsep=.6em]
\item پانکراتیت
\begin{itemize}[nosep]
\item شایع‌ترین علل \GIAarrow\ الکلیسم؛ سنگ صفراوی
\item تشخیص پانکراتیت یک تشخیص بالینی است.
\item \en{CT} برای شناسایی علت و عوارض
\end{itemize}
\item یافته‌های پانکراتیت حاد در \en{CT}
\begin{latin}\begin{itemize}[nosep]
\item Enlargement
\begin{itemize}[nosep]\item Normal size: 3 cm for head; 2.5 cm for the body; 2 cm for tail\end{itemize}
\item Peripancreatic stranding or fluid collections
\item Low-attenuation lesions from necrosis
\end{itemize}\end{latin}
\begin{itemize}[nosep]
\item نشان‌دهندهٔ پروگنوز بد
\item معمولاً تظاهر زودهنگام است.
\item نیاز به \en{CT} با کنتراست وریدی برای تشخیص
\end{itemize}
\begin{latin}\begin{itemize}[nosep]\item Pseudocyst formation\end{itemize}\end{latin}
\begin{itemize}[nosep]\item دیوارهٔ سودوسیست معمولاً با \en{CT} قابل رؤیت است و ممکن است \en{contrast-enhance} بشود.\end{itemize}
\end{enumerate}\GIAfooter{۹}

% pic10.jpg = 20261019_082732823.jpg
\clearpage\GIAheader
\begin{enumerate}[label=\arabic*.,leftmargin=*,itemsep=.6em]\setcounter{enumi}{2}
\item پانکراتیت مزمن
\begin{itemize}[nosep]
\item عموماً ناشی از سوءمصرف الکل
\item منجر به فیبروز، آتروفی، دیلاتاسیون مجرا و دیابت
\end{itemize}
\begin{latin}Diagnostic hallmark: Multiple, \textbf{amorphous calcifications} that form within the \textbf{dilated ducts} of the atrophied gland.\end{latin}
\item آدنوکارسینوم پانکراس
\begin{itemize}[nosep]
\item \en{RF}: الکل و سیگار، پانکراتیت مزمن و دیابت
\item در ۷۵\% موارد تومور در سر پانکراس است.
\item نیمی از بیماران با زردی مراجعه می‌کنند.
\item در اغلب موارد درد نیز وجود دارد هنگام مراجعه.
\end{itemize}
\textbf{* در بیمار با زردی \en{US} تصویربرداری اولیه است.}
\item یافته‌های آدنوکارسینوم پانکراس در \en{CT}
\begin{latin}\begin{itemize}[nosep]
\item Focal pancreatic mass: \textbf{Usually hypodense}
\item \textbf{Ductal dilation} $\longrightarrow$ both pancreatic and biliary ducts
\begin{itemize}[nosep]\item N/L pancreatic duct $<4$ mm in head\item CBD $<7$ mm\end{itemize}
\item Metastasis
\item Enlarged lymph nodes
\item Ascites
\end{itemize}\end{latin}
\end{enumerate}\GIAfooter{۱۰}

% pic11.jpg = 20261019_082739425.jpg
\clearpage\GIAheader
\subsection*{بیماری‌های هپاتوبیلیاری}
\begin{enumerate}[label=\arabic*.,leftmargin=*,itemsep=.6em]
\item در \en{CT} کبد \en{(Triple-phase)} \GIAarrow\ قبل و دو بار بعد از کنتراست وریدی
\begin{itemize}[nosep]\item برای بررسی توده‌های کبد\end{itemize}
\item \en{MRI} در حال تبدیل به تصویربرداری انتخابی برای بیماری‌های \en{focal} و منتشر کبد است.
\item \en{MRI} نسبت به \en{CT} در بررسی ضایعات کوچک \en{($<10$ mm)} حساس‌تر است.
\item در \en{MRI} کبد از کنتراست حاوی گادولینیوم استفاده می‌شود.
\end{enumerate}
\subsection*{انفیلتراسیون چربی}
\begin{enumerate}[label=\arabic*.,leftmargin=*,itemsep=.6em]
\item انواع انفیلتراسیون چربی
\begin{itemize}[nosep]\item منتشر\item ناحیه‌ای: متعدد؛ منفرد\end{itemize}
\item در \en{CT} بدون کنتراست کبد همواره \en{dense}تر یا با دانسیتهٔ مشابه نسبت به طحال دیده می‌شود.
\begin{itemize}[nosep]\item در انفیلتراسیون چربی کبد، طحال \en{dense}تر دیده می‌شود.\end{itemize}
\item در انفیلتراسیون منتشر کبد عموماً اندکی بزرگ می‌شود. عروق خونی نیز برجسته‌تر دیده می‌شوند.
\begin{itemize}[nosep]\item ولی انسداد یا تغییر مکان ندارند.\end{itemize}
\end{enumerate}\GIAfooter{۱۱}

% pic12.jpg = 20261019_082744787.jpg
\clearpage\GIAheader
\begin{enumerate}[label=\arabic*.,leftmargin=*,itemsep=.5em]\setcounter{enumi}{3}
\item انفیلتراسیون \en{focal} ممکن است نمای شبیه به تومور پیدا کند.
\begin{itemize}[nosep]\item اثر فشاری ندارد.\item ممکن است ظرف چند هفته پدید شده یا از بین برود.\end{itemize}
\item \en{MRI} دقیق‌ترین تصویربرداری برای بررسی کبد چرب است.
\end{enumerate}
\subsection*{سیروز}
\begin{enumerate}[label=\arabic*.,leftmargin=*,itemsep=.6em]
\item سیروز مرحلهٔ انتهایی بیماری‌های کبد است.
\begin{itemize}[nosep]\item \en{Hep B}\item \en{Hep C}\item الکلیسم\item \en{NAFLD}\item هموکروماتوز\item ویلسون\end{itemize}
\item عوارض سیروز
\begin{itemize}[nosep]\item هایپرتنشن پورت\item آسیت\item اختلال کلیوی\item \en{HCC}\item نارسایی کبد\item مرگ\end{itemize}
\end{enumerate}\GIAfooter{۱۲}

% pic13.jpg = 20261019_082752184.jpg
\clearpage\GIAheader
\begin{enumerate}[label=\arabic*.,leftmargin=*]\setcounter{enumi}{2}
\item یافته‌های سیروز در \en{CT}
\begin{latin}\begin{itemize}[nosep]
\item Diffuse fatty infiltration (Early in the disease)
\begin{itemize}[nosep]\item As the disease progresses the liver contour becomes lobulated\end{itemize}
\item The liver shrinks in volume $\longrightarrow$ Right lobe become smaller
\item Mottled \textbf{inhomogeneous} appearance to the liver parenchyma
\item Portal hypertension $\longrightarrow$ Can lead to \textbf{dilated vessels} around the stomach, splenic hilum and esophagus
\item \textbf{Splenomegaly}
\item Ascites
\end{itemize}\end{latin}
\end{enumerate}
\subsection*{ضایعات فضاگیر کبد}
\subsubsection*{متاستازها}
\begin{enumerate}[label=\arabic*.,leftmargin=*,itemsep=.5em]
\item متاستازها شایع‌ترین توده‌های بدخیم کبد هستند.
\begin{itemize}[nosep]\item اکثرشان متعدد هستند ولی شایع‌ترین بدخیمی‌های منفرد کبد نیز هستند.\end{itemize}
\item اکثر متاستازها در کبد منشأ \en{GI} دارند. \GIAarrow\ مخصوصاً کولون
\item متاستازها به کبد در \en{CT} و \en{MRI}
\begin{latin}\begin{itemize}[nosep]
\item Are usually multiple, \textbf{low attenuation} masses
\item Larger metastases may have \textbf{areas of necrosis}
\item Mucin-producing carcinomas $\longrightarrow$ from stomach, colon, ovary
\begin{itemize}[nosep]\item Can \textbf{calcify} in both primary tumors and metastases.\end{itemize}
\end{itemize}\end{latin}
\begin{itemize}[nosep]\item \en{MRI} و \en{CT} حساسیت یکسانی در تشخیص متاستاز به کبد دارند.\item \en{MRI} در نقش حل‌کنندهٔ سؤال در \en{CT}\end{itemize}
\end{enumerate}\GIAfooter{۱۳}

% pic14.jpg = 20261019_082757155.jpg
\clearpage\GIAheader
\subsection*{کارسینوم هپاتوسلولار (هپاتوما) \en{(HCC)}}
\begin{enumerate}[label=\arabic*.,leftmargin=*,itemsep=.5em]
\item شایع‌ترین بدخیمی اولیهٔ کبد است.
\item عمدتاً در کبدهای با بیماری زمینه‌ای مثل سیروز یا هپاتیت
\item \en{HCC}ها عموماً منفرد هستند \GIAarrow\ \en{$\frac{1}{5}$} متعدد \en{(DDX:} متاستاز\en{)}
\item \textbf{تهاجم عروقی در \en{HCC} شایع است.}
\begin{itemize}[nosep]\item عمدتاً به سیستم پورت\end{itemize}
\item یافته‌های \en{HCC} در \en{CT} و \en{MRI}
\begin{latin}\begin{itemize}[nosep]
\item Most HCCs are \textbf{hypodense or isodense} without contrast
\item Enhance on the arterial phase and then hypo or isodense on the venous phase
\item Low attenuation from necrosis. (Common)
\item \textbf{Calcification is frequent.}
\item In MRI
\begin{itemize}[nosep]\item Intrahepatic metastases\item Venous invasion\item HCC show washout of the contrast material.\end{itemize}
\end{itemize}\end{latin}
\end{enumerate}
\subsection*{کاورنوس همانژیوم‌ها}
\begin{enumerate}[label=\arabic*.,leftmargin=*,itemsep=.5em]
\item شایع‌ترین تومور اولیهٔ کبد است.
\item در خانم‌ها شایع‌تر است.
\end{enumerate}\GIAfooter{۱۴}

% pic15.jpg = 20261019_082805598.jpg
\clearpage\GIAheader
\begin{enumerate}[label=\arabic*.,leftmargin=*,itemsep=.5em]\setcounter{enumi}{2}
\item عموماً منفرد و تقریباً همیشه بدون علامت هستند.
\item \en{MRI} تصویربرداری انتخابی در بررسی همانژیوم‌ها است.
\item نمای کاراکتریستیک همانژیوم در \en{CT} و \en{MRI}
\begin{latin}\begin{itemize}[nosep]\item \textbf{Nodular} enhancement from the periphery inward following injection of IV contrast.\end{itemize}\end{latin}
\begin{itemize}[nosep]\item \textbf{در اسکن تأخیری \en{(10 min)} دیده می‌شود چون مادهٔ کنتراست در عروق می‌ماند.}\end{itemize}
\end{enumerate}
\begin{latin}
\textbf{Hepatic Cysts}
\begin{itemize}[nosep]
\item Congenital in origin
\item \textbf{Sharply marginated} lesions of low attenuation
\item MRI is the modality of choice
\item Solitary and homogeneous in density
\end{itemize}\end{latin}
\subsection*{سیستم صفراوی}
\begin{enumerate}[label=\arabic*.,leftmargin=*]
\item \en{MRCP} نیاز به تزریق کنتراست ندارد. طی یک دم عمیق انجام می‌شود.
\begin{itemize}[nosep]\item اگر شک به بدخیمی داشته باشیم با تزریق گادولینیوم انجام می‌دهیم.\end{itemize}
\end{enumerate}
\subsection*{مجاری ادراری}
\begin{enumerate}[label=\arabic*.,leftmargin=*]
\item کیست‌های کلیوی شایع‌ترین یافتهٔ \en{US} و \en{CT} شکم هستند.
\begin{itemize}[nosep]\item در نیمی از افراد بالای ۵۵ سال دیده می‌شوند.\end{itemize}
\end{enumerate}\GIAfooter{۱۵}

% pic16.jpg = 20261019_082810010.jpg
\clearpage\GIAheader
\begin{enumerate}[label=\arabic*.,leftmargin=*,itemsep=.5em]\setcounter{enumi}{1}
\item کیست‌های ساده
\begin{itemize}[nosep]
\item خوش‌خیم
\item \en{Fluid-filled}
\item عموماً متعدد و دوطرفه
\item هانسفیلد آب (۱۰− تا ۲۰+)
\item \textbf{مارژین شارپ در \en{CT}}
\item در پی تزریق کنتراست انهنس نمی‌کنند.
\end{itemize}
\item کیست ساده در \en{US}
\begin{latin}\begin{itemize}[nosep]\item Echo-free\item Sharp-borders\item Round or oval shape\end{itemize}\end{latin}
\end{enumerate}
\textbf{* ضخیم شدن جداره یا اکو داخلی شک به بدخیمی ایجاد می‌کند.}
\subsection*{\en{RCC} (هایپرنفروما)}
\begin{enumerate}[label=\arabic*.,leftmargin=*,itemsep=.5em]
\item شایع‌ترین بدخیمی اولیهٔ کلیه در بزرگسالان
\item توده‌های \en{solid} در کلیهٔ بزرگسالان معمولاً \en{RCC} هستند.
\item تمایل دارند وریدهای کلیوی را درگیر کنند.
\begin{itemize}[nosep]\item سپس \en{IVC} \GIAarrow\ سپس متاستاز به ریه\end{itemize}
\item وقتی به استخوان متاستاز بدهند
\begin{latin}\begin{itemize}[nosep]\item purely lytic\item often expansile\end{itemize}\end{latin}
\end{enumerate}\GIAfooter{۱۶}

% pic17.jpg = 20261019_082820220.jpg
\clearpage\GIAheader
\begin{enumerate}[label=\arabic*.,leftmargin=*,itemsep=.5em]\setcounter{enumi}{4}
\item \textbf{بررسی اولیهٔ توده‌های کلیه \GIAarrow\ \en{US}}
\item در \en{US}
\begin{itemize}[nosep]
\item \en{RCC}های کوچک \GIAarrow\ عموماً \textbf{هایپراکو} هستند.
\item \en{RCC}های بزرگ‌تر \GIAarrow\ \textbf{به‌دلیل نکروز هایپواکو می‌شوند.}
\end{itemize}
\item یافته‌های \en{RCC} در \en{CT} (باید با کنتراست باشد)
\begin{latin}\begin{itemize}[nosep]
\item Renal vein invasion (in 1/3 cases) $\longrightarrow$ filling defect in vein
\item Usually solid lesions that may contain low-attenuation areas of necrosis.
\item Still lower in density than the normal kidney (after contrast)
\end{itemize}\end{latin}
\item کاربرد \en{MRI} در کلیه: بررسی ضایعات زیر \en{1.5 cm} یا توده‌های نامشخص در \en{CT} یا \en{US}
\end{enumerate}
\subsection*{لگن}
\begin{enumerate}[label=\arabic*.,leftmargin=*,itemsep=.5em]
\item \en{US} تصویربرداری انتخابی در بررسی مشکلات لگن در زنان است.
\item \en{Transitional cell tumors}
\begin{latin}\begin{itemize}[nosep]
\item Bladder, ureter, Kidney
\item Focal thickening of the bladder wall
\item Filling defect in the contrast filled bladder
\end{itemize}\end{latin}
\end{enumerate}
\subsection*{آدنوپاتی}
\begin{enumerate}[label=\arabic*.,leftmargin=*,itemsep=.5em]
\item لنفوم می‌تواند هر نقطه‌ای را در شکم درگیر کند.
\item لنفوم‌های نودولر عموماً لنفوم هوچکین هستند.
\begin{itemize}[nosep]\item معده شایع‌ترین مکان اکسترانودال در \en{GI} برای لنفوم غیرهوچکین است.\end{itemize}
\end{enumerate}\GIAfooter{۱۷}

% pic18.jpg = 20261019_082825676.jpg
\clearpage\GIAheader
\begin{enumerate}[label=\arabic*.,leftmargin=*,itemsep=.6em]\setcounter{enumi}{3}
\item ویژگی لنفوم غیرهوچکین
\begin{itemize}[nosep]\item درگیری اکسترانودال\end{itemize}
\begin{latin}\begin{itemize}[nosep]\item Noncontiguous spread to other organs and nodes\end{itemize}\end{latin}
\item یافته‌های لنفوم شکم در \en{CT}
\begin{latin}\begin{itemize}[itemsep=.4em]
\item Multiple enlarged lymph nodes
\begin{itemize}[nosep]\item Pelvic lymph nodes are pathologic if $>1$ cm in shortest dimension\end{itemize}
\item Conglomerate masses of coalesced nodes.
\item Lymphadenopathy will displace the aorta and vena cava anteriorly.
\end{itemize}\end{latin}
\end{enumerate}\GIAfooter{۱۸}

\ifnum\GIAStandalone=1 \end{document}\fi
